% vim: set textwidth=78 autoindent:

\section{Uso dei Plugin esterni QGIS Python}\label{sec:external_plugins}\index{plugin}

% when the revision of a section has been finalized, 
% comment out the following line:
%\updatedisclaimer

I plugin esterni QGIS sono scritti in python e sono mantenuti in un archivio ufficiale moderato dall'autore. La tabella \ref{tab:external_plugins} mostra una lista di plugin attualmente disponibili con una breve descrizione.\footnote{Anche gli aggiornamenti dei plugin core possono essere disponibili in questo archivio come fonte esterna.} 
\footnote{L'installatore di Plugin Python è anch'esso un plugin Python esterno ma è parte del sorgente QGIS, caricato automaticamente e selezionabile nel gestore dei plugin QGIS Plugin Manager (vedere Sezione~\ref{sec:load_external_plugin}).}

Quando questo manuale è stato distribuito, l'archivio ufficiale, moderato e mantenuto, non era ancora completamente costituito. Una documentazione dettagliata circa l'uso, l'autore ed altre importanti informazioni sono fornite con il plugin esterno stesso e non fa parte di questo manuale.

Si può trovare una lista aggiornata dei plugin esterni moderati nell'archivio ufficiale di QGIS a \dropmenuopttwo{plugin_installer}{Scarica Python
Plugins...}) e \url{http://qgis.osgeo.org/download/plugins.html}. 

\begin{table}[H]
\centering
\caption{Attuali Plugin QGIS esterni moderati}\label{tab:external_plugins}\medskip
\small
 \begin{tabular}{|l|l|p{4in}|}
\hline \textbf{Icona} & \textbf{Plugin esterno} & \textbf{Descrizione}\\
\hline
\includegraphics[width=0.7cm]{zoom2point_icon}
 & Zoom al punto \index{plugin!zoom al punto} & Fa lo zoom alle coordinate specificate nella finestra di dialogo di input. Si può specificare il livello di zoom così come il controllo dell'estensione della vista.\\
\hline
\end{tabular}
\end{table}

Una descrizione dettagliata per i plugin python esterni si può trovare nella 
Sezione \ref{sec:load_external_plugin}.

\minisec{Archivio Plugin Python sviluppati degli utenti e Archivi degli autori}

In aggiunta ai plugin esterni moderati, esiste anche un altro archivio Plugin Python non ufficiale. Contiene plugin che non sono ancora abbastanza maturi da essere inclusi nell'archivio ufficiale, comunque alcuni possono essere abbastanza utili. Inoltre, alcuni dei nostri contributori mantengono i propri archivi.

Per aggiungere l'archivio non ufficiale e gli archivi degli autori aprire l'installatore di Plugin Python (\mainmenuopt{Plugins} > \dropmenuopttwo{plugin_installer}{Scarico Plugin Python}),
andare alla scheda \tab{Repositories} e premere \button{Aggiungi repository di terze parti}. Se non si vuole uno o più degli archivi aggiunti, disabilitarli con il tasto
\button{Edit...} o rimuoverli completamente con il tasto \button{Delete}.

\begin{Tip} \caption{\textsc{Aggiungere più Plugin esterni}}
\qgistip{In aggiunta ai plugin esterni moderati dall'archivio ufficiale QGIS si possono aggiungere più archivi esterni. Quindi selezionare la scheda Repositories nell'installatore di Plugin Pythons}
\end{Tip}

